\begin{lemma}
\label{lemma-affine-compare-bounded}
Let $X = \Spec(A)$ be an affine scheme. All the functors in the diagram
$$
\xymatrix{
D(\QCoh(\mathcal{O}_X)) \ar[rr]_{(\ref{equation-compare})}
& &
D_\QCoh(\mathcal{O}_X) \ar[ld]^{R\Gamma(X, -)} \\
& D(A) \ar[lu]^{\widetilde{\ \ }}
}
$$
are equivalences of triangulated categories. Moreover, for $E$ in
$D_\QCoh(\mathcal{O}_X)$ we have $H^0(X, E) = H^0(X, H^0(E))$.
\end{lemma}

\begin{proof}
The functor $R\Gamma(X, -)$ gives a functor
$D(\mathcal{O}_X) \to D(A)$ and hence by restriction a functor
\begin{equation}
\label{equation-back}
R\Gamma(X, -) : D_\QCoh(\mathcal{O}_X) \longrightarrow D(A).
\end{equation}
We will show this functor is quasi-inverse to (\ref{equation-compare})
via the equivalence between quasi-coherent modules on $X$ and
the category of $A$-modules.

\medskip\noindent
Elucidation. Denote $(Y, \mathcal{O}_Y)$ the one point space with sheaf
of rings given by $A$. Denote
$\pi : (X, \mathcal{O}_X) \to (Y, \mathcal{O}_Y)$
the obvious morphism of ringed spaces.
Then $R\Gamma(X, -)$ can be identified with $R\pi_*$ and the functor
(\ref{equation-compare}) via the equivalence
$\textit{Mod}(\mathcal{O}_Y) = \text{Mod}_A = \QCoh(\mathcal{O}_X)$
can be identified with $L\pi^* = \pi^* = \widetilde{\ }$ (see
Modules, Lemma \ref{modules-lemma-construct-quasi-coherent-sheaves} and
Schemes, Lemmas \ref{schemes-lemma-compare-constructions} and
\ref{schemes-lemma-equivalence-quasi-coherent}). Thus the functors
$$
\xymatrix{
D(A) \ar@<1ex>[r] & D(\mathcal{O}_X) \ar@<1ex>[l]
}
$$
are adjoint (by Cohomology, Lemma \ref{cohomology-lemma-adjoint}). In
particular we obtain canonical adjunction mappings
$$
a : \widetilde{R\Gamma(X, E)} \longrightarrow E
$$
for $E$ in $D(\mathcal{O}_X)$ and
$$
b : M^\bullet \longrightarrow R\Gamma(X, \widetilde{M^\bullet})
$$
for $M^\bullet$ a complex of $A$-modules.

\medskip\noindent
Let $E$ be an object of $D_\QCoh(\mathcal{O}_X)$. We may apply
Lemma \ref{lemma-application-nice-K-injective}
to the functor $F(-) = \Gamma(X, -)$
with $N = 1$ by Cohomology of Schemes, Lemma
\ref{coherent-lemma-quasi-coherent-affine-cohomology-zero}.
Hence
$$
H^0(R\Gamma(X, E)) = H^0(R\Gamma(X, \tau_{\geq 0}E)) = \Gamma(X, H^0(E))
$$
(the last equality by definition of the canonical truncation).
Using this we will show that the adjunction mappings $a$ and $b$
induce isomorphisms $H^0(a)$ and $H^0(b)$. Thus $a$ and $b$
are quasi-isomorphisms (as the statement is invariant under shifts)
and the lemma is proved.

\medskip\noindent
In both cases we use that $\widetilde{\ }$ is an exact functor
(Schemes, Lemma \ref{schemes-lemma-spec-sheaves}). Namely, this
implies that
$$
H^0\left(\widetilde{R\Gamma(X, E)}\right) =
\widetilde{H^0(R\Gamma(X, E))} =
\widetilde{\Gamma(X, H^0(E))}
$$
which is equal to $H^0(E)$ because $H^0(E)$ is quasi-coherent. Thus
$H^0(a)$ is an isomorphism. For the other direction we have
$$
H^0(R\Gamma(X, \widetilde{M^\bullet})) =
\Gamma(X, H^0(\widetilde{M^\bullet})) =
\Gamma(X, \widetilde{H^0(M^\bullet)}) =
H^0(M^\bullet)
$$
which proves that $H^0(b)$ is an isomorphism.
\end{proof}
